\chapter{对偶群与 Langlands 函子性}\label{ch:11}

\begin{center}
J.\,W. Cogdell
\end{center}

Langlands 从未把 $\GL_n$ 的 Langlands 猜想与他的一般函子性原则 [30] 分开。
特别地，对任意连通约化代数群 $G$，他表述了某些 Galois 表示与容许表示
或自守表示之间的一个对应。对 $\GL_n$，这是某些 $n$ 维 Galois 表示
（即映入 $\GL_n(\mathbb{C})$ 的表示）与 $\GL_n(k)$ 的容许表示
或 $\GL_n(\mathbb{A})$ 的自守表示之间的对应 [4]。对一般的 $G$，
自守侧该换成什么是清楚的：$G(k)$ 的容许表示，或 $G(\mathbb{A})$ 的自守表示。
但 Galois 侧的目标 $\GL_n(\mathbb{C})$ 该换成什么？

基于无分歧表示的 Satake 参数化 [33]，以及他对代数环面表示的分类 [24]，
Langlands 引入了对偶群的想法，今称朗兰兹对偶群或 L 群，记作 ${}^{L}G$，
用以充当 $\GL_n(\mathbb{C})$ 的角色。$n$ 维 Galois 表示的角色，
则由 Galois 群映入这个 L 群的某些容许同态承担。就函子性而言，
最方便的看法是：把 $G$ 的这些局部与整体对应，看成用这些映入 ${}^{L}G$
的容许 Galois 同态，给出 $G$ 的容许表示或自守表示的算术参数化。
\begin{annote}
「算术参数化」是本章的枢纽。它的意思是：不再从表示出发去造 Galois 表示
（那是非 Abel 类域论的方向），而是反过来把 Galois 侧的数据当作标签，
贴到自守侧的表示上。这样函子性就退化为「复合两个标签映射」这一句话，
见 \S3 的讨论。须注意标签只在 L 包的层面唯一，单个表示一般无法由参数
唯一确定（$\GL_n$ 是例外，L 包为单点）。
\par\vspace{0.15em}{\raggedleft 【来源：Langlands, \textit{Problems in the theory of automorphic forms}, Lecture Notes in Math.\ 170 (1970), 18--86；Borel, \textit{Automorphic $L$-functions}, Proc.\ Symp.\ Pure Math.\ 33 (1979), Part 2, 27--61（A）；本译本第 10 章 \S3（I）】\par}
\end{annote}

Langlands 的函子性原则断言：任一 L 同态 ${}^{L}H\to{}^{L}G$ 都应决定一个
从 $H$ 的容许表示或自守表示到 $G$ 的容许表示或自守表示的
\emph{转移}或\emph{提升}。一旦有了参数化，这在概念上便水到渠成：
把 $H$ 的表示的参数化同态与这个 L 同态复合，即得 $G$ 的某个表示的
参数化同态。若取 $H=\{1\}$，则 ${}^{L}H$ 就是 Galois 群或与之密切相关的群，
于是本质上从这条函子性原则即可收回 $G$ 的局部或整体 Langlands 对应。

用迹公式方法已确立了许多基本的函子性例子：循环基变换、循环自守诱导、
内型之间的提升。然而近来，用 $\GL_n$ 的逆定理得到的、到 $\GL_n$ 的整体函子性
有很大进展。其中包括 Ramakrishnan [31] 从 $\GL_2\times\GL_2$ 到 $\GL_4$
的张量积提升，Kim 与 Shahidi [21] 从 $\GL_2\times\GL_3$ 到 $\GL_6$ 的同类提升，
Kim 与 Shahidi [20, 21, 22] 从 $\GL_2$ 到 $\GL_4$ 与 $\GL_5$ 的对称三次幂
与对称四次幂提升，以及 Kim、Piatetski-Shapiro 与 Shahidi [5, 6]
从分裂典型群到 $\GL_N$ 的提升。

本文先描述 L 群的构造，以及一般约化群 $G$ 的局部与整体 Langlands 猜想的
表述 [2]。接着概述 Langlands 的函子性原则及其与局部、整体 Langlands 对应的
关系。然后转向例子：先简述上述用迹公式确立的若干函子性例子，
再较详细地描述到 $\GL_n$ 的新函子性，以及如何用逆定理作为确立这些提升的手段。

感谢审稿人帮助澄清与本章有关的若干问题。

\section{对偶群}

设 $G$ 为定义在 $k$ 上的连通约化代数群，$k$ 为局部域或整体域。
令 $\bar k$ 为 $k$ 的可分代数闭包，$\mathcal{G}_k=\operatorname{Gal}(\bar k/k)$
为 Galois 群。在 $\bar k$ 上 $G$ 变为分裂的，并由其根数据分类 [2, 36]。
在 $G_{/\bar k}$ 中取 Borel 子群 $B$ 与极大环面 $T$，二者都在 $\bar k$ 上有定义
且分裂。令 $X=\mathcal{X}^{*}(T)$ 记 $T$ 的 $\bar k$ 有理特征之集，
$\Phi=\Phi(G,T)\subset X$ 为联系于 $G$ 与 $T$ 的根系，$\Delta\subset\Phi$
为对应于 $B$ 的单根之集。与三元组 $(X,\Phi,\Delta)$ 对偶地，
我们有三元组 $(X^{\vee},\Phi^{\vee},\Delta^{\vee})$，由余特征格
$X^{\vee}=\mathcal{X}_{*}(T)$（即 $\bar k$ 有理的单参数子群）、
余根系 $\Phi^{\vee}$ 与单余根 $\Delta^{\vee}$ 组成。四元组
$\Psi(G)=(X,\Phi,X^{\vee},\Phi^{\vee})$ 称为 $G$ 在 $\bar k$ 上的根数据，
四元组 $\Psi_0(G)=(X,\Delta,X^{\vee},\Delta^{\vee})$ 称为 $G$ 在 $\bar k$ 上的
带基根数据 [2, 36]。连通约化 $\bar k$ 群的基本结构定理如下 [36]。

\par\medskip\noindent\textbf{定理 1.1.}\quad
根数据 $\Psi(G)$ 在不计 $\bar k$-同构的意义下决定 $G$。

关于相对结构理论，有分裂正合列
\[
  1\longrightarrow\operatorname{Int}(G)\longrightarrow\operatorname{Aut}(G)
  \longrightarrow\operatorname{Aut}(\Psi_0(G))\longrightarrow 1.
\]
取定每个 $\alpha\in\Delta$ 的根向量 $x_\alpha$，即给出 $G$ 的一个分裂
$(G,B,T,\{x_\alpha\}_{\alpha\in\Delta})$，并给出典范同构
\[
  \operatorname{Aut}(\Psi_0(G))\longrightarrow
  \operatorname{Aut}(G,B,T,\{x_\alpha\})\subset\operatorname{Aut}(G).
\]
\begin{annote}
这里的「分裂」是英文 splitting 的双关：既指正合列的分裂（同态
$\operatorname{Aut}(\Psi_0(G))\to\operatorname{Aut}(G)$），又指把 $G$ 连同
$(B,T,\{x_\alpha\})$ 钉住的一套数据（今多称 pinning）。之所以钉住根向量就能
唯一确定自同构，是因为单根处的根向量一旦给定，整个根系上的根向量由
$\mathrm{ad}$-作用与 Chevalley 交换关系唯一恢复，故自同构在
$\mathfrak{g}$ 上的作用被完全锁定。定理 1.1 本身是 Chevalley 的存在与唯一性
定理：给定根数据必有群，且在同构意义下唯一，见 Springer 的教材。
\par\vspace{0.15em}{\raggedleft 【来源：Chevalley, \textit{Classification des groupes de Lie alg\'ebriques} (1956--58)；Springer, \textit{Linear Algebraic Groups}, 2nd ed., Progress in Math.\ 9, Birkh\"auser 1998, \S\,14--15（A）；Borel--Tits, Publ.\ Math.\ IHES 27 (1965), \S\,1--2（A）】\par}
\end{annote}

若 $G$ 定义在 $k$ 上，则 $\mathcal{G}_k$ 作用在 $G_{/\bar k}$ 上，给出 $k$-结构。
于是有同态
\[
  \mathcal{G}_k\longrightarrow\operatorname{Aut}(G_{/\bar k})
  \longrightarrow\operatorname{Aut}(\Psi_0(G)).
\]
故 $G_{/k}$ 决定两部分数据：决定 $\bar k$ 上群者，即根数据 $\Psi(G)$；
以及同态 $\mathcal{G}_k\to\operatorname{Aut}(\Psi_0(G))$。

定义 ${}^{L}G$ 只需把这套结构理论对偶化。令
$\Psi_0(G)^{\vee}=(X^{\vee},\Delta^{\vee},X,\Delta)$ 为对偶带基根数据。
它在 $\mathbb{C}$ 上定义一个连通约化代数群 ${}^{L}G^{0}$。Galois 结构可以搬过来，
因为
\[
  \operatorname{Aut}\bigl(\Psi_0({}^{L}G^{0})\bigr)
  =\operatorname{Aut}\bigl(\Psi_0(G)^{\vee}\bigr)
  =\operatorname{Aut}\bigl(\Psi_0(G)\bigr),
\]
而 ${}^{L}G^{0}$ 的上述正合列的一个分裂给出映射
$\mu\colon\mathcal{G}_k\to\operatorname{Aut}(\Psi_0({}^{L}G^{0}))
\to\operatorname{Aut}({}^{L}G^{0})$。此映射固定 ${}^{L}G^{0}$ 的相应分裂
$({}^{L}G^{0},{}^{L}B^{0},{}^{L}T^{0},\{x_{\alpha^{\vee}}\}_{\alpha^{\vee}\in\Delta^{\vee}})$，
从而给出复约化群 ${}^{L}G^{0}$ 上的 $\mathcal{G}_k$ 作用，
它编码了 $G$ 原有的部分 $k$-结构。

\par\medskip\noindent\textbf{定义 1.1.}\quad
$G$ 的（Langlands）对偶群，或 L 群，是
\[
  {}^{L}(G_{/k})={}^{L}G={}^{L}G^{0}\rtimes\mathcal{G}_k.
\]
\begin{annote}
原书此处印作 ${}^{L}G^{0}\times\mathcal{G}_k$，但 $\mathcal{G}_k$ 经 $\mu$
非平凡地作用在 ${}^{L}G^{0}$ 上，所以这是半直积，译文改排 $\rtimes$。
要紧的是：L 群不是抽象群，而是带 $\mathcal{G}_k$ 作用的复约化群；
忘掉这个作用，注 2 与条件 (iii) 的「相干」就都无从谈起。另须分清三种形式：
Galois 形式 ${}^{L}G^{0}\rtimes\mathcal{G}_k$、Weil 形式
${}^{L}G^{0}\rtimes W_k$、Weil--Deligne 形式 ${}^{L}G^{0}\rtimes W'_k$。
原书 §4 讨论 $H=\{1\}$ 时用的是第三种，因为此时 ${}^{L}H=W'_k$ 才能
与容许同态的定义域对上。
\par\vspace{0.15em}{\raggedleft 【来源：Borel, Proc.\ Symp.\ Pure Math.\ 33 (1979), \S\,2--3（A）；nLab, \textit{Langlands dual group}（B，与 Borel 交叉核对）】\par}
\end{annote}

\par\medskip\noindent\textbf{注.}
1. 有时方便使用对偶群的 Weil 形式。由于有自然映射 $W_k\to\mathcal{G}_k$，
可以改取 ${}^{L}G={}^{L}G^{0}\rtimes W_k$，二者没有本质差别。
对某些目的也可用 Weil--Deligne 形式。

2. 若 $G'$ 是 $k$ 群且在 $\bar k$ 上与 $G$ 同构，则 $G$ 与 $G'$ 互为内型，
当且仅当 ${}^{L}G$ 与 ${}^{L}G'$ 在 $\mathcal{G}_k$ 上同构 [2]。
故对偶群并不能完全区分 $k$-形式，它只区分到内型；但它完全决定准分裂形式。
\begin{annote}
机制如下。$k$-形式由 Galois 上同调 $H^1(k,G^{\mathrm{ad}})$ 分类，
而 $\mathcal{G}_k$ 对 ${}^{L}G^{0}$ 的作用只经过商
$\operatorname{Out}(G)=\operatorname{Aut}(\Psi_0(G))$，即只记录上闭链在
外自同构群中的像。内型恰是像相同的那些类，故 ${}^{L}G$ 分辨不出它们；
准分裂型对应像为平凡的那一类，故被完全决定。反过来说，
「${}^{L}G$ 相同」并不等于「$G$ 相同」，本章后面 $\Phi(G)$ 与 $\Phi(G')$
可以不等，正是这一点的表现。
\par\vspace{0.15em}{\raggedleft 【来源：Borel, Proc.\ Symp.\ Pure Math.\ 33 (1979), \S\,2；Kottwitz, \textit{Stable trace formula: cuspidal tempered terms}, Duke Math.\ J.\ 51 (1984), \S\,1--2（A）】\par}
\end{annote}

实践中，这种对偶保持 $A_n$ 与 $D_n$ 型，交换 $B_n$ 与 $C_n$ 型。
此外它还交换相应群的伴随型与单连通型。
\begin{center}
\begin{tabular}{cc}
\hline
$G$ & ${}^{L}G^{0}$ \\
\hline
$\GL_n$ & $\GL_n(\mathbb{C})$ \\
$\mathrm{SO}_{2n+1}$ & $\mathrm{Sp}_{2n}(\mathbb{C})$ \\
$\mathrm{Sp}_{2n}$ & $\mathrm{SO}_{2n+1}(\mathbb{C})$ \\
$\mathrm{SO}_{2n}$ & $\mathrm{SO}_{2n}(\mathbb{C})$ \\
伴随型 & 单连通型 \\
单连通型 & 伴随型 \\
\hline
\end{tabular}
\end{center}
\begin{annote}
表的最后两行是通例：对偶把特征格与余特征格互换，故单连通型（$X^{\vee}$
生成整个 $X^{\vee}$）与伴随型（$X$ 生成整个 $X$）互换。$\mathrm{SO}_{2n}$
一行值得停一下：$D_n$ 型有三个 $k$-形式，即单连通的 $\mathrm{Spin}_{2n}$、
伴随的 $\mathrm{PSO}_{2n}$，以及中间的 $\mathrm{SO}_{2n}$（$\mathrm{Spin}_{2n}$
模中心的对角 $\mu_2$）。$\mathrm{SO}_{2n}$ 既非单连通也非伴随，
故不在「互换」之列，其对偶仍是自身，原书直接把它列入表中而不加说明。
$B_n$ 与 $C_n$ 互换则是 $\mathrm{SO}_{2n+1}\leftrightarrow\mathrm{Sp}_{2n}$
这一经典事实。
\par\vspace{0.15em}{\raggedleft 【来源：Borel, Proc.\ Symp.\ Pure Math.\ 33 (1979), \S\,3；Springer, \textit{Linear Algebraic Groups}, 2nd ed., \S\,17（A）】\par}
\end{annote}

局部与整体的构造是相容的。故若 $G$ 定义在整体域 $k$ 上，$v$ 为 $k$ 的一个位点，
并令 $G_v$ 记 $G$ 作为 $k_v$ 上的群，则有自然映射 ${}^{L}G_v\to{}^{L}G$。

\section{\texorpdfstring{$G$}{G} 的 Langlands 猜想}

\subsection{局部 Langlands 猜想}

设 $k$ 为局部域，$W'_k$ 为相应的 Weil--Deligne 群 [4]。若 $k$ 为阿基米德域，
就简单地取 $W'_k=W_k$ 为 Weil 群。仿照 Borel [2]，称同态
$\phi\colon W'_k\to{}^{L}G$ 为\emph{容许的}，如果满足以下各条：

\begin{enumerate}[label=(\roman*),leftmargin=*]
  \item $\phi$ 是 $\mathcal{G}_k$ 上的同态，即下图交换：
\[
\begin{tikzcd}
  W'_k \arrow[r,"\phi"] \arrow[d] & {}^{L}G \arrow[d] \\
  \mathcal{G}_k \arrow[r] & \mathcal{G}_k
\end{tikzcd}
\]
  \item $\phi$ 连续，$\phi(\mathbb{G}_a)$ 在 ${}^{L}G^{0}$ 中幂幺，
    且 $\phi$ 把半单元映为半单元；
  \item 若 $\phi(W'_k)$ 含于 ${}^{L}G$ 的某个真抛物子群 $P$ 的 Levi 子群之中，
    则 $P$ 是相干的。
\end{enumerate}

对所有未定义的概念（如相干），我们请读者参阅 Borel [2]。
若 $G=\GL_n$，容许同态恰是 $W'_k$ 的 Frobenius 半单复表示 [4]。
\begin{annote}
两处须留心。其一，字母 $\Phi$ 在本章承担两种含义：\S1 的
$\Phi=\Phi(G,T)$ 是根系，本节起 $\Phi(G)$ 是容许同态之集，原书自己在
p.~254 提醒勿与文献 [4] 中表示几何 Frobenius 的 $\Phi$ 相混。读公式时
只能靠上标与括号区分。其二，\emph{相干}（relevant）原书不定义，
只指向 Borel [2]：粗略地说，$P$ 的 Levi 因子须与 $G$ 的 $k$-结构相容，
即其 ${}^{L}G^{0}$ 部分来自 $G$ 的某个 $k$-定义 Levi 子群。
条件 (iii) 排除的正是那些「在 $\bar k$ 上出现、但在 $k$ 上没有对应物」的 Levi，
这也解释了为何内型可以改变 $\Phi(G)$。此外，(iii) 与 (iv) 是全章反复使用的
判别准则：离散谱系列对应「参数不落在真 Levi 中」，缓增对应「参数像有界」。
但两者的定义域不同：(iii) 用 $W'_k$，(iv) 用 $W_k$。缓增性只看半单部分，
而 $\mathbb{G}_a$ 的像必幂幺，故原书 p.~255 的排法并非笔误，不可随手统一。
\par\vspace{0.15em}{\raggedleft 【来源：Borel, Proc.\ Symp.\ Pure Math.\ 33 (1979), \S\,6--8（A）；Langlands, \textit{On the classification of irreducible representations of real algebraic groups}（A）；本译本第 10 章 \S3（I，术语「容许同态」）】\par}
\end{annote}

仿照 Borel [2] 与 Langlands [26]，以 $\Phi(G)$ 记所有容许同态
$\phi\colon W'_k\to{}^{L}G$ 模去 ${}^{L}G^{0}$ 中元素的内自同构之集。
注意若 $G$ 与 $G'$ 互为内型，从而 ${}^{L}G={}^{L}G'$，
则未必有 $\Phi(G)=\Phi(G')$，因为上述条件 (iii) 看得见 $k$-结构。
若 $G$ 是准分裂形式，则确实有 $\Phi(G')\subset\Phi(G)$。

为陈述 $G$ 的局部 Langlands 猜想，还需两个补充构造，请参阅 Borel [2]。
首先，对每个 $\phi\in\Phi(G)$，有一种方式构造 $G$ 的中心 $C(G)$ 的特征
$\omega_\phi$。其次，若以 $C({}^{L}G^{0})$ 记 ${}^{L}G^{0}$ 的中心，
则对每个 $\alpha\in H^{1}\bigl(W'_k;C({}^{L}G^{0})\bigr)$ 都联系着 $G(k)$
的一个特征 $\chi_\alpha$。若把 $\phi\in\Phi(G)$ 写成
$\phi=(\phi_1,\phi_2)$，其中 $\phi_1(w)\in{}^{L}G^{0}$、
$\phi_2(w)\in\mathcal{G}_k$，则 $\phi_1$ 是 $W'_k$ 上取值于 ${}^{L}G^{0}$
的闭链，映射 $\phi\mapsto\phi_1$ 给出嵌入
$\Phi(G)\hookrightarrow H^{1}(W'_k;{}^{L}G^{0})$。于是
$H^{1}\bigl(W'_k;C({}^{L}G^{0})\bigr)$ 自然作用在
$H^{1}(W'_k;{}^{L}G^{0})$ 上，且此作用保持 $\Phi(G)$。

有了这些构造，便可陈述 $G$ 的局部 Langlands 猜想 [2]。如前，
以 $\mathcal{A}(G)=\mathcal{A}(G(k))$ 记 $G(k)$ 的不可约容许复表示的
等价类之集。

\par\medskip\noindent\textbf{局部 Langlands 猜想.}\quad
设 $k$ 为局部域。则有满射 $\mathcal{A}(G)\to\Phi(G)$，纤维有限，
它把 $\mathcal{A}(G)$ 分成互不相交的有限集
$\mathcal{A}_\phi=\mathcal{A}_\phi(G)$，满足以下各条：
\begin{enumerate}[label=(\roman*),leftmargin=*]
  \item 若 $\pi\in\mathcal{A}_\phi$，则 $\pi$ 的中心特征 $\omega_\pi$
    等于 $\omega_\phi$；
  \item 与扭相容：若 $\alpha\in H^{1}\bigl(W'_k;C({}^{L}G^{0})\bigr)$，
    $\chi_\alpha$ 为相应的 $G(k)$ 的特征，则
    $\mathcal{A}_{\alpha\cdot\phi}=\{\pi\chi_\alpha\mid\pi\in\mathcal{A}_\phi\}$；
  \item 一个元 $\pi\in\mathcal{A}_\phi$ 模 $C(G)$ 平方可积，
    当且仅当所有 $\pi\in\mathcal{A}_\phi$ 模 $C(G)$ 平方可积，
    当且仅当 $\phi(W'_k)$ 不含于 ${}^{L}G$ 的真 Levi 子群中；
  \item 一个元 $\pi\in\mathcal{A}_\phi$ 缓增，当且仅当所有
    $\pi\in\mathcal{A}_\phi$ 缓增，当且仅当 $\phi(W_k)$ 有界；
  \item 若 $H$ 为连通约化 $k$ 群，$\eta\colon H(k)\to G(k)$ 为核与余核
    都交换的 $k$ 态射，则 $G(k)$ 与 $H(k)$ 的分解之间应有相容性。
    即：$\eta$ 诱导自然映射 ${}^{L}\eta\colon{}^{L}G\to{}^{L}H$；
    对 $\phi\in\Phi(G)$ 置 $\phi'={}^{L}\eta\circ\phi$，
    则任一 $\pi\in\mathcal{A}_\phi(G)$ 视为 $H(k)$ 模时，
    分解为有限多个属于 $\mathcal{A}_{\phi'}(H)$ 的不可约容许表示的直和。
\end{enumerate}

对 $\phi\in\Phi(G)$，集合 $\mathcal{A}_\phi(G)$ 称为 \emph{L 包}。
文献 [4] 中局部 Langlands 猜想的版本 I 是此猜想对群 $\GL_n$ 的特殊化。
在那种情形，所有 L 包都是单点集，$\mathcal{A}(G)\to\Phi(G)$ 是双射。
此猜想给出 $G(k)$ 的不可约容许表示的算术参数化。

除 $\GL_n$ 的结果外，关于此猜想已知以下情形。
\begin{enumerate}[label=\arabic*.,leftmargin=*]
  \item 若局部域 $k$ 阿基米德，即 $k=\mathbb{R}$ 或 $\mathbb{C}$，
    则由 Langlands [26] 完全建立。
  \item 若 $k$ 非阿基米德，$G$ 在 $k$ 上准分裂，且在某个有限 Galois 扩张上分裂，
    则已知如何用无分歧容许同态参数化 $G(k)$ 的无分歧表示 [2]。
    这是用这套语言重述 Satake 分类 [33]。
  \item 若 $k$ 非阿基米德，则 Kazhdan 与 Lusztig 已说明如何用 Weil--Deligne 群的
    容许同态，参数化 $G(k)$ 中有 Iwahori 不动向量的那些表示 [19]。
  \item 近来，在 $k$ 为非阿基米德特征零、$G$ 为分裂 $\mathrm{SO}_{2n+1}$ 的情形，
    Jiang 与 Soudry 用 Weil--Deligne 群的容许同态给出了
    $\mathrm{SO}_{2n+1}(k)$ 的 generic 表示的参数化 [16, 17]。
    他们由 $\GL_{2n}(k)$ 的参数化拉回而得，是近期关于从分裂
    $\mathrm{SO}_{2n+1}$ 到 $\GL_{2n}$ 的整体函子性工作（后文将讨论）的产物。
\end{enumerate}
若把此版本的局部 Langlands 猜想视为给出 $G(k)$ 不可约容许表示的算术参数化，
则可用此参数化对任意 $\pi\in\mathcal{A}(G)$ 定义局部 L 函数。这需要第二个参数，
即表示 $r\colon{}^{L}G\to\GL_n(\mathbb{C})$：这里指连续同态，
且其在 ${}^{L}G^{0}$ 上的限制是复 Lie 群态射。于是对任意容许同态
$\phi\in\Phi(G)$，复合 $r\circ\phi\colon W'_k\to\GL_n(\mathbb{C})$
是 [4] 所考虑的 Weil--Deligne 群的连续复表示，可对其联系 L 因子
$L(s,r\circ\phi)$ 与 $\varepsilon$ 因子 $\varepsilon(s,r\circ\phi,\psi)$，
其中 $\psi$ 为 $k$ 的加法特征。

\par\medskip\noindent\textbf{定义 2.1.}\quad
若 $\pi\in\mathcal{A}_\phi$ 位于容许同态 $\phi$ 所定义的 L 包中，则置
\[
  L(s,\pi,r)=L(s,r\circ\phi),\qquad
  \varepsilon(s,\pi,r,\psi)=\varepsilon(s,r\circ\phi,\psi).
\]
\begin{annote}
这个定义把 L 函数的归一化完全交给 $r\circ\phi$，而 $\phi$ 的定义域是
Weil--Deligne 群，故与第 10 章所用的一致：那里的几何 Frobenius 使
$\omega^{s}(\Phi)=q^{-s}$，此处 $L(s,r\circ\phi)$ 取 [4] 的约定，
两者不差 $q$ 的幂。若换用算术 Frobenius，所有 $\varepsilon$ 因子须整体扭一个
$|\cdot|^{s}$ 型的特征，且 $\varepsilon(s,\pi,r,\psi)$ 的 $\psi$ 也要相应
乘一个幂。对 $\GL_n$，此定义与 Rankin--Selberg 积分给出的解析定义一致，
这正是版本 II 的内容；对一般 $G$ 则只有参数化已知时才能这样定义。
\par\vspace{0.15em}{\raggedleft 【来源：Cogdell, 本书第 9 章；Cogdell, 本书第 10 章 \S1--2（I，归一化约定）；Tate, \textit{Number theoretic background}, Proc.\ Symp.\ Pure Math.\ 33 (1979), Part 2, \S\,3--4（A）】\par}
\end{annote}

按此定义，同一 L 包 $\mathcal{A}_\phi$ 中的表示 $\pi$ 无法用其 L 函数
与 $\varepsilon$ 因子区分，这正是「L 包」这一术语的由来。目前，
这些 L 函数只对参数化已知的那些 $\pi$ 有良好定义，例如 $\pi$ 无分歧时。

若以此作为容许表示所附局部 L 函数的定义，则文献 [4] 中局部 Langlands 猜想的
版本 II 会表述为：以解析方式定义的匹配的 L 因子与 $\varepsilon$ 因子
（如 [3] 中对 $\GL_n$ 者）与此处定义者相等。对一般约化群，
我未见按此方式的表述；但在上引 Jiang 与 Soudry 的工作中，他们做到的正是这一点。
他们对 $\mathrm{SO}_{2n+1}(k)$ 的每个 generic 表示 $\pi$ 附上一个容许同态
$\phi_\pi$，使得对标准嵌入
$r\colon\mathrm{Sp}_{2n}(\mathbb{C})\hookrightarrow\GL_{2n}(\mathbb{C})$ 有等式
\[
\begin{aligned}
  L(s,\pi\times\pi')
  &=L(s,\pi\times\pi',r\otimes\mathrm{id})\\
  &=L(s,(r\circ\phi_\pi)\otimes\rho_{\pi'}),
\end{aligned}
\]
以及相应的 $\varepsilon$ 因子等式。这里 $\pi'$ 为 $\GL_m(k)$ 的不可约容许表示，
$\rho_{\pi'}$ 为 $\GL_m$ 的局部 Langlands 猜想所附的 $W'_k$ 的相应表示，
$L(s,\pi\times\pi')$ 为 Shahidi [34] 定义的解析 L 函数。

\subsection{整体 Langlands 猜想}

现取 $k$ 为整体域，$\mathbb{A}$ 为其阿代尔环。对 $k$ 上的约化代数群 $G$，
以 $\mathcal{A}(G)=\mathcal{A}(G(\mathbb{A}))$ 记 $G(\mathbb{A})$ 的
不可约自守表示之集。与 $\GL_n$ 一样，要表述整体 Langlands 猜想，
我们须把 Weil--Deligne 群换成猜想中的 Langlands 群 $\mathcal{L}_k$，
并考虑容许同态 $\phi\colon\mathcal{L}_k\to{}^{L}G$ 之集。
这些同态应当以某种方式参数化 $G(\mathbb{A})$ 的不可约自守表示。
其确切形式目前还相当猜想性。

虽然不知道它该是什么样子，人们仍期待整体与局部相容。若从 $G(\mathbb{A})$
的不可约自守表示 $\pi=\bigotimes'_v\pi_v$ 出发，并假设每个局部群 $G(k_v)$
的局部 Langlands 猜想，则可对 $\pi$ 附上局部参数
$\phi_v=\phi_{\pi_v}\colon W'_{k_v}\to{}^{L}G_v$ 的汇集 $\{\phi_v\}$，
由局部分量 $\pi_v$ 给出。把这些与对偶群的自然相容映射
$\iota_v\colon{}^{L}G_v\to{}^{L}G$ 复合，即得局部参数
$\iota_v\circ\phi_v\colon W'_{k_v}\to{}^{L}G$ 的汇集
$\{\iota_v\circ\phi_v\}$。

这样的映射系统必须来自整体参数 $\phi\colon\mathcal{L}_k\to{}^{L}G$，
局部与整体理论才自洽。这个局部参数系统常常可以充当整体参数 $\phi$ 的替代物。
例如，这组局部数据足以定义联系于 $\pi$ 的整体 L 函数与 $\varepsilon$ 因子。
若 $r\colon{}^{L}G\to\GL_n(\mathbb{C})$，则复合
$r_v=r\circ\iota_v\colon{}^{L}G_v\to\GL_n(\mathbb{C})$ 给出局部对偶群的表示。

\par\medskip\noindent\textbf{定义 2.2.}\quad
若 $\pi=\bigotimes'_v\pi_v$ 为 $G(\mathbb{A})$ 的不可约自守表示，
$r\colon{}^{L}G\to\GL_n(\mathbb{C})$，则置
\[
  L(s,\pi,r)=\prod_v L(s,\pi_v,r_v)
  =\prod_v L(s,r\circ\iota_v\circ\phi_v),
\]
以及
\[
  \varepsilon(s,\pi,r)=\prod_v\varepsilon(s,\pi_v,r_v,\psi_v)
  =\prod_v\varepsilon(s,r\circ\iota_v\circ\phi_v,\psi_v),
\]
其中 $\psi=\bigotimes_v\psi_v$ 为 $\mathbb{A}$ 的、在 $k$ 上平凡的加法特征。
\begin{annote}
两个无穷乘积都有意义，因为对几乎所有的 $v$，$r_v$ 无分歧，
$\varepsilon(s,\pi_v,r_v,\psi_v)=1$，而 $L$ 因子是 $q_v^{-s}$ 的多项式的倒数，
次数由 $\dim r$ 界定。故真正的解析内容在于：这些局部因子的乘积能否延拓、
满足函数方程。原书随即指出，完整的整体 L 函数需要所有位点的局部 Langlands 猜想，
这只对 $\GL_n$ 成立；一般 $G$ 只能先取定有限集 $S$ 而谈部分 L 函数。
部分 L 函数与完整者的差别只在有限多个因子，故函数方程只差一个显式的
有限乘积，这也是逆定理可以只在部分 L 函数上操作的原因。
\par\vspace{0.15em}{\raggedleft 【来源：Langlands, \textit{Euler Products}, Yale Univ.\ Press 1971（A）；Borel, Proc.\ Symp.\ Pure Math.\ 33 (1979), \S\,10（A）】\par}
\end{annote}

如上述定义完整 L 函数，须在所有位点解出局部 Langlands 猜想，
而这仅对 $\GL_n$ 已知。不过，对任一不可约自守表示 $\pi$，存在有限位点集
$S=S(\pi)$，使得对所有 $v\notin S$，表示 $\pi_v$ 无分歧，
从而局部参数化问题已解决。于是部分 L 函数
\[
  L^{S}(s,\pi,r)=\prod_{v\notin S}L(s,\pi_v,r_v)
\]
总有良好定义，且 Langlands 已证明此 Euler 乘积在某个右半平面绝对收敛 [25]。

\section{函子性}

从其近作 [29, 30] 可以看出，Langlands 一向把「函子性原则」
视为他关于自守表示的看法之核心。它把上文所称的「局部与整体 Langlands 猜想」
作为特例包含在内。

令 $k$ 为局部域或整体域，$H$ 与 $G$ 为 $k$ 上两个连通约化群。
我们已定义它们相应的对偶群 ${}^{L}H$ 与 ${}^{L}G$。称同态
$u\colon{}^{L}H\to{}^{L}G$ 为 \emph{L 同态}，如果
(i) 它是 $\mathcal{G}_k$ 上的同态，即下图交换：
\[
\begin{tikzcd}
  {}^{L}H \arrow[r,"u"] \arrow[d] & {}^{L}G \arrow[d] \\
  \mathcal{G}_k \arrow[r] & \mathcal{G}_k
\end{tikzcd}
\]
(ii) $u$ 连续；(iii) $u$ 在 ${}^{L}H^{0}$ 上的限制是复解析同态
$u\colon{}^{L}H^{0}\to{}^{L}G^{0}$。

此外若 $G$ 准分裂，则对任意容许同态 $\phi\in\Phi(H)$，
复合 $u\circ\phi$ 仍是 $\Phi(G)$ 中的容许同态。于是映射
$\phi\mapsto u\circ\phi$ 定义了 $\Phi(u)\colon\Phi(H)\to\Phi(G)$。
若 $k$ 为整体域、$v$ 为 $k$ 的位点，则由 $\mathcal{G}_{k_v}$ 可自然视为
$\mathcal{G}_k$ 的子群，可把 ${}^{L}G_v$ 视为 ${}^{L}G$ 的子群。
于是限制到 ${}^{L}H_v$ 上，$u$ 诱导局部对偶群的 L 同态
$u_v\colon{}^{L}H_v\to{}^{L}G_v$，从而有局部映射
$\Phi(u_v)\colon\Phi(H_v)\to\Phi(G_v)$。

函子性原则现在可粗略表述如下 [30]。

\par\medskip\noindent\textbf{函子性原则.}\quad
若 $k$ 为局部（相应整体）域，$H$ 与 $G$ 为连通约化 $k$ 群且 $G$ 准分裂，
则对每个 L 同态 $u\colon{}^{L}H\to{}^{L}G$，都联系着一个从 $H$ 的
容许（相应自守）表示到 $G$ 的容许（相应自守）表示的转移或提升。
\begin{annote}
「$G$ 准分裂」这一条不可省。容许同态与条件 (iii) 的相干性都是相对
准分裂形式定义的（见 \S2 的讨论：内型共用同一个 ${}^{L}G$，
但 $\Phi(G')$ 只是 $\Phi(G)$ 的子集）。若 $G$ 非准分裂，
$u\circ\phi$ 可能落在 $\Phi(G)$ 之外，提升就无从谈起。原书随后也只对
准分裂的 $G$ 陈述定义 3.1。实践中的目标群——$\GL_n$、典型群的分裂形式——
都满足此条，所以不妨碍应用。
\par\vspace{0.15em}{\raggedleft 【来源：Langlands, \textit{Problems in the theory of automorphic forms} (1970), \S\,2--3；Borel, Proc.\ Symp.\ Pure Math.\ 33 (1979), \S\,9（A）】\par}
\end{annote}

若假设局部与整体 Langlands 猜想，即已有 $\mathcal{A}(H)$ 与 $\mathcal{A}(G)$
的算术参数化，则这个提升过程容易描述。

\subsection{局部函子性}

先取 $k$ 为局部域，$u\colon{}^{L}H\to{}^{L}G$ 为局部 L 同态。
若取 $\pi\in\mathcal{A}(H)$ 为 $H(k)$ 的不可约容许表示，
则它由容许同态 $\phi=\phi_\pi\colon W'_k\to{}^{L}H$ 参数化。
事实上，$\phi$ 参数化整个局部 L 包 $\mathcal{A}_\phi(H)$。
把 $\phi$ 与 $u$ 复合，得 $\phi'=\Phi(u)(\phi)=u\circ\phi\in\Phi(G)$，
即 $W'_k$ 到 ${}^{L}G$ 的容许同态。于是 $\phi'$ 参数化局部 L 包
$\mathcal{A}_{\phi'}(G)$，而这个 L 包（有时是其中任一元素 $\Pi$）
就是 $\pi$ 或 L 包 $\mathcal{A}_\phi(H)$ 的函子性提升
（或转移、Langlands 提升等）。

一般地，在我们理解局部参数化的情形，也就「理解」了局部函子性：
\begin{enumerate}[label=\arabic*.,leftmargin=*]
  \item $k=\mathbb{R}$ 或 $\mathbb{C}$，$H$ 为任意连通约化 $k$ 群，
    $G$ 为任意准分裂连通约化 $k$ 群；
  \item $k$ 为非阿基米德局部域，$H=\GL_m$，$G=\GL_n$
    （及相关例子，见第 4 节）；
  \item 设 $k$ 非阿基米德，整数环为 $\mathcal{O}$。设 $H$ 与 $G$ 都准分裂，
    且存在 $k$ 的有限扩张 $K$，使 $H$ 与 $G$ 都在 $K$ 上分裂并有 $\mathcal{O}$
    结构，使 $H(\mathcal{O})$ 与 $G(\mathcal{O})$ 都是特殊极大紧子群。
    设 $\pi$ 为 $H(k)$ 的无分歧表示，具有非平凡的 $H(\mathcal{O})$ 不动向量，
    无分歧参数 $\phi=\phi_\pi\in\Phi(H)$。则对任意 L 同态
    $u\colon{}^{L}H\to{}^{L}G$，参数 $\phi'=u\circ\phi$ 无分歧，
    并定义 L 包 $\mathcal{A}_{\phi'}(G)$，其中含（唯一的）$G(k)$ 的表示
    $\Pi$，它关于 $G(\mathcal{O})$ 无分歧 [2]。$\Pi$ 称为 $\pi$ 的
    \emph{自然无分歧提升}。
\end{enumerate}

\subsection{整体函子性}

若现在考虑 $k$ 为整体域，则原则上函子性提升应如局部情形一样，
用整体参数化来运作。但现在我们再次处于不利地位：并不真正理解参数化群
$\mathcal{L}_k$。作为替代，我们回落到所期望的局部--整体相容性。
故设 $H$ 为连通约化 $k$ 群，$G$ 为准分裂连通约化 $k$ 群，
$u\colon{}^{L}H\to{}^{L}G$ 为 L 同态。对 $k$ 的每个位点 $v$，
有上述相应的局部 L 同态 $u_v\colon{}^{L}H_v\to{}^{L}G_v$。
现设 $\pi\in\mathcal{A}(H)$、$\pi=\bigotimes'_v\pi_v$ 为 $H(\mathbb{A})$
的不可约自守表示。若 $v$ 阿基米德，则由 Langlands 的工作我们知道如何用
局部参数 $\phi_v\colon W'_{k_v}\to{}^{L}H_v$ 参数化 $\pi_v$。
若 $v$ 非阿基米德，则对几乎所有的 $v$，局部群 $H_v$ 准分裂、
在 $k_v$ 的某有限扩张上分裂，且 $\pi_v$ 关于某特殊极大紧子群无分歧。
故我们处于已有 $\pi_v$ 的局部参数
$\phi_v\colon W'_{k_v}\to{}^{L}H_v$ 的情形。在这两种情形，
都可构造局部的提升 $\Pi_v$，即联系于参数 $\phi'_v=u_v\circ\phi_v$ 的
$G(k_v)$ 的表示，也就是上述定义的局部提升。

\par\medskip\noindent\textbf{定义 3.1.}\quad
设 $H$ 为连通约化 $k$ 群，$\pi=\bigotimes'_v\pi_v$ 为 $H(\mathbb{A})$
的不可约自守表示。设 $G$ 为准分裂连通约化 $k$ 群，
$u\colon{}^{L}H\to{}^{L}G$ 为 L 同态。则 $G(\mathbb{A})$ 的自守表示
$\Pi=\bigotimes'_v\Pi_v$ 称为 $\pi$ 的（弱）函子性提升（关于 $u$），
如果对所有阿基米德位点，以及 $\pi_v$ 无分歧的几乎所有有限位点，
$\Pi_v$ 都是上述关于 $u_v$ 的局部函子性提升。若 $\Pi_v$ 对 $k$ 的所有位点
都是 $\pi_v$ 的局部函子性提升，则称 $\Pi$ 为 $\pi$ 的（强）函子性提升。
\begin{annote}
原书此段之前有一处笔误：先说「现设 $\pi\in\mathcal{A}(G)$，
$\pi=\bigotimes'\pi_v$ 为 $H(\mathbb{A})$ 的不可约自守表示」，
按上下文应为 $\mathcal{A}(H)$，译文已改正。定义中「弱」与「强」的差别
只在有限多个位点，看似无关紧要，实则关键：弱提升的存在性可由逆定理得到，
而它的唯一性来自 $\GL_N$ 的强重数一（两个自守表示若几乎处处相同，
则整体相同），所以弱提升一旦存在就是确定的。$\mathrm{SO}_{2n+1}$ 情形
之所以能由此反推局部参数，靠的正是这一点。
\par\vspace{0.15em}{\raggedleft 【来源：Jacquet--Shalika, \textit{On Euler products and the classification of automorphic representations I}, Amer.\ J.\ Math.\ 103 (1981), 499--558（A）；Cogdell--Piatetski-Shapiro, \textit{Converse theorems for $\GL_n$ II}, J.\ reine angew.\ Math.\ 507 (1999)（A）】\par}
\end{annote}

注意，由此定义可得：若 $\pi$ 为 $H(\mathbb{A})$ 的自守表示，
$u\colon{}^{L}H\to{}^{L}G$ 为 L 同态，$\Pi$ 为 $\pi$ 到 $G(\mathbb{A})$
的自守表示的函子性提升，则对每个表示 $r\colon{}^{L}G\to\GL_n(\mathbb{C})$，
都有 L 函数与 $\varepsilon$ 因子的等式
\[
  L^{S}(s,\pi,r\circ u)=L^{S}(s,\Pi,r),\qquad
  \varepsilon^{S}(s,\pi,r\circ u,\psi)=\varepsilon^{S}(s,\Pi,r,\psi),
\]
其中 $S$ 为我们不知道如何局部提升 $\pi_v$ 的有限（可能空）位点集。

事实上，我们需要在 L 包的层面上做这件事。这很容易表述，
但鉴于我们知识的残缺，此时这样做似乎收获甚微。不过，由局部与整体 L 包现象
带来的局部提升、从而整体提升的歧义，应始终牢记。

\section{例子}

我们已指出，Langlands 视函子性为包含局部与整体 Langlands 猜想及其推论
（如强 Artin 猜想）者。原因之一便是下面这个例子。

考虑 $H=\{1\}$ 的情形。先取 $k$ 为局部域。由于有从 Weil--Deligne 群
$W'_k$ 到 $\mathcal{G}_k$ 的自然映射，可考虑 L 群的 Weil--Deligne 形式
${}^{L}G={}^{L}G^{0}\rtimes W'_k$。于是 ${}^{L}H=W'_k$。例如取
$G=\GL_n$，则 $u\colon{}^{L}H\to{}^{L}G$ 就是 $\Phi(G)$ 中的容许同态，
即 Weil--Deligne 群的复表示，而这些群的函子性包含局部 Langlands 猜想。
若取 $k$ 为整体域，并让 ${}^{L}G$ 保持 L 群的 Galois 形式，
则同样取 $H=\{1\}$、$G=\GL_n$，即得 $\GL_n$ 的整体 Langlands 猜想。

我愿讨论的其他函子性例子，可分为我看来两类：Galois 理论的与群论的。
前者包括基变换、自守诱导与内型之间的提升。后者都是到 $\GL_n$ 的提升，
包括张量积提升、对称幂提升与来自典型群的提升。我不会触及称为
\emph{内窥}（endoscopic）的重要一类提升，尽管所讨论的某些例子可解释为
（可能扭的）内窥的例子。内窥提升是那些 L 同态 $u\colon{}^{L}H\to{}^{L}G$
把 ${}^{L}H^{0}$ 实现为 ${}^{L}G^{0}$ 中某对合的不动点的提升。
这类提升的重要性主要在于理解迹公式时不可或缺，而此处无法讨论。
请读者参阅 Langlands [28, 30] 以及 Kottwitz 与 Shelstad [23] 的工作
与其中文献。

\subsection{Galois 理论的例子}

在这些例子中，L 同态源自 Galois 理论。

\par\medskip\noindent\textbf{1. 基变换（或自守限制）.}\quad
设 $K$ 为 $k$ 的有限扩张。则在 Weil 群层面有 $W_K\subset W_k$，
故 $W_k$ 的任一表示经限制给出 $W_K$ 的表示。容许表示或自守表示层面的
相应提升如下。设 $H$ 连通、约化且在 $k$ 上分裂。则可把 $H$ 也视为 $K$ 上的群；
令 $G=R_{K/k}(H)$ 为 Weil 的从 $K$ 到 $k$ 的纯量限制，故 $G(k)=H(K)$，
而 $G$ 是由 $H_{/K}$ 决定的 $k$ 上的群。于是有自然嵌入
\[
\begin{aligned}
  u\colon{}^{L}H&={}^{L}H^{0}\times\mathcal{G}_k\\
  &\longrightarrow\Bigl(\prod_{\mathcal{G}_K\backslash\mathcal{G}_k}
  {}^{L}H^{0}\Bigr)\rtimes\mathcal{G}_k={}^{L}G,
\end{aligned}
\]
其中 $\mathcal{G}_k$ 通过置换指标集作用在 $\prod{}^{L}H^{0}$ 上，
而在 ${}^{L}H^{0}$ 上是对角映射、在 $\mathcal{G}_k$ 上是恒等。
当 $k$ 为局部域时，诱导映射 $\Phi(u)\colon\Phi(H)\to\Phi(G)$ 确实是限制映射，
把 $W'_K$ 视为 $W'_k$ 的开子群。来自此 L 同态的函子性，
会从 $H(k)$ 或 $H(\mathbb{A}_k)$ 的表示 $\pi$ 出发，
产生 $G(k)=H(K)$ 或 $H(\mathbb{A}_K)$ 的表示，称为 $\pi$ 的\emph{基变换}。
此纲领在 $H=\GL_n$、扩张 $K/k$ 可解时已完成：$n=2$ 由 Langlands [27]，
一般 $n$ 由 Arthur 与 Clozel [1]。其技巧是扭迹公式。此外，$H=\GL_2$ 时
Jacquet、Piatetski-Shapiro 与 Shalika 用逆定理方法得到了非正规的三次基变换 [14]。
\begin{annote}
指标集 $\mathcal{G}_K\backslash\mathcal{G}_k$ 是双陪集之集，
与 $k$-嵌入 $K\hookrightarrow\bar k$ 的集合 $\operatorname{Hom}_k(K,\bar k)$
一一对应，故当 $K/k$ 为 $d$ 次扩张时它有 $d$ 个元素；
$\mathcal{G}_k$ 在其上的作用就是 Galois 作用，置换这些嵌入。
这正是纯量限制 $R_{K/k}H$ 的对偶群含 $d$ 个 ${}^{L}H^{0}$ 因子的原因，
也是 $R_{K/k}H$ 在 $\bar k$ 上同构于 $H^{d}$ 的原因。
对角嵌入则保证：限制到 $W_K$ 后，$d$ 个因子取同一个表示，
这正是「限制」在 L 群层面的对应物。
\par\vspace{0.15em}{\raggedleft 【来源：Arthur--Clozel, \textit{Simple Algebras, Base Change, and the Advanced Theory of the Trace Formula}, Ann.\ of Math.\ Studies 120, Princeton 1989, Ch.\ 1（A）；Borel, Proc.\ Symp.\ Pure Math.\ 33 (1979), \S\,5（A）】\par}
\end{annote}

\par\medskip\noindent\textbf{2. 自守诱导.}\quad
仍取 $K$ 为 $k$ 的 $d$ 次有限可分扩张，故 $W_K\subset W_k$。
若从 $W_K$ 的表示出发，只需归纳即得 $W_k$ 的表示。
容许或自守表示层面的相应提升如下。取 $H=R_{K/k}(\GL_n)$，
即把 $\GL_n(K)$ 视为上述 $k$ 上的群，并令 $G=\GL_{dn}(k)$。
现在有 L 同态
\[
\begin{aligned}
  u\colon{}^{L}H&=\Bigl(\prod_{\mathcal{G}_K\backslash\mathcal{G}_k}
  \GL_n(\mathbb{C})\Bigr)\rtimes\mathcal{G}_k\\
  &\longrightarrow{}^{L}G=\GL_{dn}(\mathbb{C})\rtimes\mathcal{G}_k,
\end{aligned}
\]
把 ${}^{L}H^{0}=\GL_n(\mathbb{C})\times\cdots\times\GL_n(\mathbb{C})$
作为分块对角矩阵映入 ${}^{L}G^{0}=\GL_{dn}(\mathbb{C})$，
并让 $\mathcal{G}_k$ 通过 $\mathfrak{S}_d$ 中的置换矩阵作用在
$\GL_{dn}(\mathbb{C})$ 上，从而延拓为 L 同态。来自这种 L 同态的局部或整体
函子性称为\emph{自守诱导}。容许同态集上的映射 $\Phi(u)$ 应当是归纳。
同样，当扩张 $K/k$ 可解时，Arthur 与 Clozel [1] 用扭迹公式在局部与整体上
作了分析，$n=2$ 的先驱是 Jacquet 与 Langlands [15]。Henniart 与 Herb 在
Kazhdan 关于 $n=1$ 情形的早期工作 [18] 基础上，用局部特征恒等式首次给出
并分析了 $\GL_n$ 的局部自守诱导 [13]。此工作所用的迹公式版本比 [1] 或 [18]
都简单，且容许正特征域。
\begin{annote}
为何用 $\mathfrak{S}_d$ 的置换矩阵？因为 ${}^{L}G^{0}=\GL_{dn}(\mathbb{C})$
中分块对角矩阵的正规化子模去对角块，恰是 $\mathfrak{S}_d$（Weyl 群的一部分），
而 $\mathcal{G}_k$ 的像必须落在其中，才能与分块结构相容。
维数上也可核对：$\operatorname{Ind}_{W_K}^{W_k}\rho$ 的维数是
$d\cdot\dim\rho$，与分块对角的 $d$ 个 $\GL_n(\mathbb{C})$ 块相符。
故 $\Phi(u)$ 取归纳而非限制，与基变换恰成对照：
一个是把 $W_k$ 的表示限制到 $W_K$，一个是把 $W_K$ 的表示归纳到 $W_k$。
\par\vspace{0.15em}{\raggedleft 【来源：Arthur--Clozel, Ann.\ of Math.\ Studies 120, Ch.\ 3（A）；Henniart--Herb, \textit{Automorphic induction for $GL(n)$}, Duke Math.\ J.\ 78 (1995), 131--192（A）】\par}
\end{annote}

\par\medskip\noindent\textbf{3. 内型.}\quad
设 $G$ 在局部或整体 $k$ 上连通、约化、准分裂，$H$ 为 $G$ 的内型。
则 ${}^{L}H={}^{L}G$，恒等映射 $u\colon{}^{L}H\to{}^{L}G$ 是 L 同态，
故应有相应的提升。注意若 $k$ 为局部域，则 $\Phi(H)\subset\Phi(G)$；
若 $k$ 为整体域，则事实上对几乎所有位点 $H_v=G_v$，
故 $\Phi(H_v)=\Phi(G_v)$。在 $G=\GL_2$、$H=D^{\times}$
（$k$ 上秩 2 除法代数的乘法群）的情形，从 $D^{\times}$ 的表示到 $\GL_2$
的表示的这个提升，就是所谓的 Jacquet--Langlands 对应，已在 [15] 中建立。
若取 $G=\GL_n$、$H=\GL_m(D)$，其中 $D$ 是秩 $d$ 的中心单除法代数、
$dm=n$，则局部函子性已由 Rogawski [32] 在 $m=1$ 情形、由 Deligne、Kazhdan
与 Vignéras [10] 用迹公式作了分析。

\subsection{群论的例子}

在这组例子中，涉及的群 $H$ 都分裂，目标群 $G$ 总是通用线性群 $\GL_n$，
故 Galois 理论作用不大。L 同态是来自群论的自然映射。近来，
以 $\GL_n$ 的逆定理作为确立到 $\GL_n$ 的整体函子性的主要工具，
这一族函子性有很大进展。

\par\medskip\noindent\textbf{1. 张量积.}\quad
设 $k$ 为局部或整体域，$H=\GL_m\times\GL_n$。则
${}^{L}H^{0}=\GL_m(\mathbb{C})\times\GL_n(\mathbb{C})$，
${}^{L}H={}^{L}H^{0}\times\mathcal{G}_k$。若取 $G=\GL_{mn}$，
则 ${}^{L}G^{0}=\GL_{mn}(\mathbb{C})$，${}^{L}G={}^{L}G^{0}\times\mathcal{G}_k$。
简单的张量积映射
$\otimes\colon\GL_m(\mathbb{C})\times\GL_n(\mathbb{C})\to\GL_{mn}(\mathbb{C})$
以 $\mathcal{G}_k$ 上的恒等映射延拓，定义了 L 同态
$u_\otimes\colon{}^{L}H\to{}^{L}G$。相应的函子性是\emph{张量积提升}。
注意若 $k$ 为局部域，则局部提升原则上现已理解，
因为 $\GL_n$ 的局部参数化问题（局部 Langlands 猜想）已解决。
故有趣的问题是整体函子性。近来这已在 $\GL_2\times\GL_2$ 到 $\GL_4$
（Ramakrishnan [31]）与 $\GL_2\times\GL_3$ 到 $\GL_6$
（Kim 与 Shahidi [21]）的情形解决。

\par\medskip\noindent\textbf{2. 对称幂.}\quad
设 $k$ 为局部或整体域，$H=\GL_2$，故 ${}^{L}H^{0}=\GL_2(\mathbb{C})$，
${}^{L}H={}^{L}H^{0}\times\mathcal{G}_k$。取 $G=\GL_{n+1}$（$n\ge1$），
故 ${}^{L}G^{0}=\GL_{n+1}(\mathbb{C})$，${}^{L}G={}^{L}G^{0}\times\mathcal{G}_k$。
对每个 $n\ge1$ 有自然的 $n$ 次对称幂映射
$\mathrm{sym}^{n}\colon\GL_2(\mathbb{C})\to\GL_{n+1}(\mathbb{C})$。
把它以 Galois 群上的恒等映射延拓，得 L 同态
$\mathrm{sym}^{n}\colon{}^{L}H\to{}^{L}G$。相应的函子性是从 $\GL_2$
的表示到 $\GL_{n+1}$ 的表示的\emph{对称幂提升}。同样，若 $k$ 为局部域，
则因 $\GL_n$ 的局部 Langlands 猜想已解决，局部对称幂提升原则上已理解。
故有趣的仍是整体函子性。整体对称平方提升（即 $\GL_2$ 到 $\GL_3$）
是 Gelbart 与 Jacquet 的旧定理 [11]。近来，Kim 与 Shahidi 证明了从 $\GL_2$
到 $\GL_4$ 的整体对称三次幂提升的存在性 [21, 22]，随后 Kim 又给出从 $\GL_2$
到 $\GL_5$ 的整体对称四次幂提升 [20, 22]。对所有 $n$ 实现对称幂函子性，
将导致 $\GL_2$ 的 Ramanujan 猜想的证明。
\begin{annote}
这里给出的是经典论证的轮廓。设 $\pi$ 对应未分歧参数
$\mathrm{diag}(\alpha,\beta)$，则 $\mathrm{sym}^{n}\pi$ 的参数为
$\alpha^{n},\alpha^{n-1}\beta,\dots,\beta^{n}$。若对所有 $n$ 都有对称幂提升，
则每个 $L(s,\mathrm{sym}^{n}\pi)$ 都是自守 L 函数，从而在 $\Re s>1$
绝对收敛（自守 L 函数在收敛半平面内无零点极点之外的问题）。
把 Euler 因子展开，得 $\sum_{n}|{\alpha}|^{n}q^{-ns}$ 型级数的收敛性，
取 $n\to\infty$ 便得 $|\alpha|\le1$，同理 $|\beta|\le1$；
结合中心特征给出的 $|\alpha\beta|=1$，即得 $|\alpha|=|\beta|=1$，
这正是 Ramanujan 猜想。故对称幂函子性是通往 Ramanujan 的最短路线。
\par\vspace{0.15em}{\raggedleft 【来源：Langlands, \textit{Problems in the theory of automorphic forms} (1970), \S\,5；Shahidi, \textit{Symmetric power $L$-functions}（A）】\par}
\end{annote}

\par\medskip\noindent\textbf{3. 典型群.}\quad
仍取 $k$ 为局部或整体域。取 $H$ 为 $k$ 上的分裂典型群，
更具体地，是 $\mathrm{SO}_{2n+1}$、$\mathrm{Sp}_{2n}$ 或 $\mathrm{SO}_{2n}$
的分裂形式。则 L 群的连通分量分别是 $\mathrm{Sp}_{2n}(\mathbb{C})$、
$\mathrm{SO}_{2n+1}(\mathbb{C})$ 或 $\mathrm{SO}_{2n}(\mathbb{C})$，
且都有到适当通用线性群的自然嵌入。
\begin{center}
\begin{tabular}{ccccc}
\hline
$H$ & ${}^{L}H^{0}$ & $u^{0}\colon{}^{L}H^{0}\hookrightarrow{}^{L}G^{0}$
  & ${}^{L}G^{0}$ & $G$ \\
\hline
$\mathrm{SO}_{2n+1}$ & $\mathrm{Sp}_{2n}(\mathbb{C})$
  & $\mathrm{Sp}_{2n}(\mathbb{C})\hookrightarrow\GL_{2n}(\mathbb{C})$
  & $\GL_{2n}(\mathbb{C})$ & $\GL_{2n}$ \\
$\mathrm{Sp}_{2n}$ & $\mathrm{SO}_{2n+1}(\mathbb{C})$
  & $\mathrm{SO}_{2n+1}(\mathbb{C})\hookrightarrow\GL_{2n+1}(\mathbb{C})$
  & $\GL_{2n+1}(\mathbb{C})$ & $\GL_{2n+1}$ \\
$\mathrm{SO}_{2n}$ & $\mathrm{SO}_{2n}(\mathbb{C})$
  & $\mathrm{SO}_{2n}(\mathbb{C})\hookrightarrow\GL_{2n}(\mathbb{C})$
  & $\GL_{2n}(\mathbb{C})$ & $\GL_{2n}$ \\
\hline
\end{tabular}
\end{center}
\begin{annote}
三行嵌入都是「标准表示」：$\mathrm{Sp}_{2n}$ 的标准表示是 $2n$ 维，
$\mathrm{SO}_{2n+1}$ 的是 $2n+1$ 维，$\mathrm{SO}_{2n}$ 的是 $2n$ 维。
复合 $r=\iota\circ u^{0}$ 正是这些标准表示，故提升后的 L 函数
$L(s,\pi,r)$ 就是该典型群表示的「标准 L 函数」。
表中 $H$ 与 ${}^{L}H^{0}$ 的对应已在 \S1 的表中出现过，
这里再列一次是为了与 $u^{0}$ 及目标 $G$ 对齐：一旦选定嵌入，
提升的目标群与 L 函数就同时确定。
\par\vspace{0.15em}{\raggedleft 【来源：Borel, Proc.\ Symp.\ Pure Math.\ 33 (1979), \S\,3--4；Cogdell--Kim--Piatetski-Shapiro--Shahidi, Publ.\ Math.\ IHES 93 (2001), \S\,1（A）】\par}
\end{annote}

这些同态以 Galois 群上的恒等映射延拓，即得 L 同态。与每个相应，
应有从 $\mathcal{A}(H)$ 到 $\mathcal{A}(G)$ 的容许或自守表示的提升。
与 Kim、Piatetski-Shapiro 和 Shahidi 合作，我们用逆定理方法对数域 $k$ 上
从 $\mathrm{SO}_{2n+1}$ 到 $\GL_{2n}$ 的 generic 尖点表示建立了整体弱提升 [5]。
此后不久，Ginzburg、Rallis 与 Soudry 证明我们的弱提升确实是强提升，
并刻画了其像 [12]。上引 Jiang 与 Soudry 关于 $p$ 进域上
$\mathrm{SO}_{2n+1}$ 的局部 Langlands 猜想的结果 [16, 17]，
即作为此整体函子性的局部推论得到。我们也已能把函子性结果推广到
其他分裂典型群 [6]。

我们想解释用逆定理获得到通用线性群的整体函子性的方法。从数域 $k$ 上
定义的群 $H$ 出发。取 $\pi=\bigotimes_v\pi_v$ 为 $H(\mathbb{A})$ 的尖点表示。
对每个局部位点 $v$，施用局部函子性，构造 $G(k_v)=\GL_N(k_v)$ 的局部表示
$\Pi_v$（$N$ 适当）。若属上例 1 或 2，则可对所有 $v$ 这样做，
因为 $\GL_n(k_v)$ 的局部 Langlands 猜想已知 [4]。对典型群的情形，
可在所有阿基米德位点 $v$ 以及 $\pi_v$ 无分歧的非阿基米德位点 $v$ 这样做。
方法不过是用上述 L 同态复合 $\pi_v$ 的局部参数映射 $\phi_v$。
典型群情形，须在其余位点 $v$ 巧取局部提升以构造局部提升 $\Pi_v$。
但暂设我们理解所有位点的局部提升。则由构造有局部 L 因子的等式
\[
  L(s,\pi_v,r_v)=L(s,r_v\circ\phi_v)=L(s,u_v\circ\phi_v)
  =L(s,\Pi_v,\iota_v),
\]
以及局部 $\varepsilon$ 因子的类似等式。这里可取 $r=u^{0}$，
视为复表示 $r\colon{}^{L}H\to\GL_N(\mathbb{C})$，而
$\iota\colon{}^{L}G\to\GL_N(\mathbb{C})$ 不过是到第一个因子 ${}^{L}G^{0}$
的投影。于是若置 $\Pi=\bigotimes'_v\Pi_v$ 为 $\GL_N(\mathbb{A})$ 的
不可约容许表示，则整体有
\[
  L(s,\pi,r)=L(s,\Pi,\iota),\qquad
  \varepsilon(s,\pi,r)=\varepsilon(s,\Pi,\iota).
\]
此外，若 $\pi'=\bigotimes_v\pi'_v$ 为 $\GL_m(\mathbb{A})$ 的尖点表示，
$m\le N-2$，则同样有
\[
  L(s,\pi_v\times\pi'_v,r_v\otimes\iota_v)
  =L(s,\Pi_v\times\pi'_v,\iota_v\otimes\iota_v),
\]
从而
\[
\begin{aligned}
  L(s,\pi\times\pi',r\otimes\iota)
  &=L(s,\Pi\times\pi',\iota\otimes\iota)\\
  &=L(s,\Pi\times\pi'),
\end{aligned}
\]
以及局部与整体 $\varepsilon$ 因子的类似等式。如 [3] 所述，要应用 $\GL_N$
的逆定理，须对足够多的尖点扭 $\pi'$ 控制扭 L 函数
$L(s,\Pi\times\pi')=L(s,\Pi\times\pi',\iota\otimes\iota)$ 的解析性质。
但由 L 因子与 $\varepsilon$ 因子的等式，这些性质由群 $H(\mathbb{A})$ 的自守
L 函数 $L(s,\pi\times\pi',r\otimes\iota)$ 在 $\GL_m(\mathbb{A})$ 扭下的
解析性质所控制。故一旦已知这些 L 函数的足够解析控制，只需应用 $\GL_N$
的逆定理 [3]，即得 $\Pi$ 自守。至今多数情形，这种对
$L(s,\pi\times\pi',r\otimes\iota)$ 的解析控制，都是由 Langlands--Shahidi
方法、通过 Eisenstein 级数的 Fourier 系数分析 L 函数而得到的。

现按此梗概重访上述例子。

\par\medskip\noindent\textbf{1. 张量积.}\quad
Ramakrishnan [31] 的情形，函子性从 $\GL_2\times\GL_2$ 到 $\GL_4$，
$\pi=\pi_1\otimes\pi_2$，各 $\pi_i$ 为 $\GL_2(\mathbb{A})$ 的尖点表示，
$\Pi$ 为 $\GL_4(\mathbb{A})$ 的自守表示。要应用 [9] 的逆定理，
Ramakrishnan 须对 $\GL_1(\mathbb{A})$ 与 $\GL_2(\mathbb{A})$ 的尖点表示
$\pi'$ 控制 $L(s,\pi\times\pi')$ 的解析性质，即 Rankin 三重积 L 函数
\[
  L(s,\pi\times\pi',r\otimes\iota)=L(s,\pi_1\times\pi_2\times\pi').
\]
他得以用 Garrett 及随后 Rallis 与 Piatetski-Shapiro 关于此 L 函数的
积分表示，加上 Shahidi 关于 Langlands--Shahidi 方法的工作，完成此事。
Kim 与 Shahidi [21] 的情形类似，只是 $\pi_2$ 现为 $\GL_3(\mathbb{A})$
的尖点表示。不过，由于提升得到的表示 $\Pi$ 是 $\GL_6(\mathbb{A})$
的自守表示，要应用 [9] 的逆定理，他们须控制
\[
  L(s,\Pi\times\pi')=L(s,\pi_1\times\pi_2\times\pi')
\]
的解析性质，其中 $\pi'$ 须跑遍 $\GL_m(\mathbb{A})$ 的适当尖点表示，
$m=1,2,3,4$。这些三重积的控制是 Langlands--Shahidi 方法的应用，
涉及 $\GL_5$、$\mathrm{Spin}_{10}$ 以及单连通的 $E_6$ 与 $E_7$ 上
Eisenstein 级数的系数。

\par\medskip\noindent\textbf{2. 对称幂.}\quad
Gelbart 与 Jacquet 最初的对称平方提升确实用了 $\GL_3$ 的逆定理 [11]。
只需控制特征的扭（即 $\GL_1$ 上的自守形式），须控制的 L 函数是 $\GL_2$
的对称平方 L 函数，因为
\[
  L(s,\Pi)=L(s,\pi,\mathrm{sym}^{2}).
\]
他们借助 Shimura 的积分表示完成了此事。对 Kim 与 Shahidi，
对称三次幂与四次幂提升是由上述函子性的 $\GL_2\times\GL_3$ 张量积提升
与 $\GL_4$ 的外平方提升推出的 [20, 21, 22]。

\par\medskip\noindent\textbf{3. 典型群.}\quad
这里有个次生问题。若从 $H(\mathbb{A})$ 的 generic 尖点表示
$\pi=\bigotimes_v\pi_v$ 出发，则有有限个有限位点组成的集合 $S$，
在其上不知道 $\pi_v$ 用容许同态表示的局部参数化，
从而不知道正确的局部提升 $\Pi_v$ 该是什么。此时可在这些位点取任意的
局部提升，只要它有平凡中心特征。为补偿，采用 $\GL_N$ 逆定理的如下形式：
固定单个高度分歧的伊代尔类特征 $\eta$，其分歧程度依赖于原有的
$H(\mathbb{A})$ 的表示 $\pi$ 与所构造的 $\GL_N(\mathbb{A})$ 的表示 $\Pi$
（实际上只依赖于 $v\in S$ 处的局部分量），然后用一切形如
$\pi'=\tau\otimes\eta$ 的 $\GL_m(\mathbb{A})$ 的尖点表示（$m\le N-1$，
$\tau$ 在 $S$ 中所有 $v$ 处无分歧）作扭 [3, 5]。这个高度分歧的扭起两个作用。
第一，它帮助控制 $H(\mathbb{A})$ 的扭 L 函数 $L(s,\pi\times\pi')$
的整体极点；第二，它使我们能通过局部 $\gamma$ 因子在高度分歧扭下的稳定性，
匹配那些 $v\in S$ 处的局部 L 因子与 $\varepsilon$ 因子 [3, 8, 5]。
故对这些受限的扭，确实有
\[
\begin{aligned}
  L(s,\pi\times\pi')&=L(s,\pi\times\pi',r\otimes\iota)\\
  &=L(s,\Pi\times\pi',\iota\otimes\iota)=L(s,\Pi\times\pi'),
\end{aligned}
\]
以及 $\varepsilon$ 因子的类似等式。由于对此族 $\pi'$ 能用 Langlands--Shahidi
方法控制 $L(s,\pi\times\pi')$ 的解析性质，我们可应用 $\GL_N$ 的逆定理，
得出存在 $\GL_N(\mathbb{A})$ 的自守表示 $\Pi'$，使对一切 $v\notin S$ 有
$\Pi_v=\Pi'_v$，即 $\pi$ 的弱提升 $\Pi'$。
\begin{annote}
这里用的是稳定性的如下形式：固定维数，当特征 $\chi$ 充分分歧时，
局部 $\gamma$ 因子 $\gamma(s,\pi_v\times\chi,\psi)$ 与 $\pi_v$ 无关，
只依赖 $\chi$ 与 $\psi$。于是即便 $\pi_v$ 本身未知，只要把两侧都扭上同一个
高度分歧的 $\eta$，两侧的 $\gamma$ 因子就自动相等，逆定理所需的局部匹配
便可强行成立——这是「以扭代知」的典型手法，代价是只能对特殊形状的
$\pi'=\tau\otimes\eta$ 使用结论。原书 p.~265 的论证依赖的正是这一点。
\par\vspace{0.15em}{\raggedleft 【来源：Cogdell--Piatetski-Shapiro, \textit{Stability of gamma factors for $SO(2n+1)$}, Manuscripta Math.\ 95 (1998), 437--461（A）；Shahidi, \textit{Local coefficients as Mellin transforms of Bessel functions}, IMRN 39 (2002), 2075--2119（A）】\par}
\end{annote}

此论证的每一步现在对所考虑的一般分裂典型群都成立。最初，
局部 $\gamma$ 因子的稳定性只对 $\mathrm{SO}_{2n+1}$ 已知 [8, 5]。如今，
多亏 Shahidi 把其局部系数表为 Bessel 函数的 Mellin 变换的近期结果 [35]，
可用 [8] 的技巧对其他分裂典型群也建立局部 $\gamma$ 因子的稳定性。
这便允许把 [5] 的函子性结果推广到这些情形 [6]。

在 $\mathrm{SO}_{2n+1}$ 情形，一旦有了弱提升，Ginzburg、Rallis 与 Soudry
的理论 [12] 便允许证明此弱提升确实是强提升，意即那些 $v\in S$ 处的
局部分量被完全决定——事实上没有可能的歧义。与此相结合，
他们得以完全刻画其像。一旦知道这些提升是刚性的，便可把 $\pi_v$ 的提升取为
整体决定的 $\Pi_v$，从而开始定义局部提升。这就是 Jiang 与 Soudry 的论文
[16, 17] 的内容。本质上，他们证明此局部提升满足人们从函子性所期待的
L 函数关系，然后用此提升拉回局部 $\GL_N(k_v)$ 表示的参数
$\phi_{\Pi_v}$（由局部 Langlands 猜想已知其存在），得到正确类型的参数
$\phi_{\pi_v}\colon W'_{k_v}\to{}^{L}H_v$，从而推出 $H(k_v)$ 的
局部 Langlands 猜想。更多细节与精确陈述请参阅他们的论文。
我们预期其他分裂典型群也会有类似结果。

\par\medskip
\noindent\textit{（原书本章末列参考文献 36 条，见原书 pp.~265--268，即底本 PDF 267--270，
编号 [1]--[36] 与正文引用一一对应；条目为书目信息，本译本不逐条迻录，
读者可按编号核对原书。）}
